
\documentclass{article}
\usepackage{colortbl}
\usepackage{makecell}
\usepackage{multirow}
\usepackage{supertabular}

\begin{document}

\newcounter{utterance}

\centering \large Interaction Transcript for game `cladder', experiment `full\_v1.5\_default', episode 9577 with llama{-}3.1{-}8b{-}instruct.
\vspace{24pt}

{ \footnotesize  \setcounter{utterance}{1}
\setlength{\tabcolsep}{0pt}
\begin{supertabular}{c@{$\;$}|p{.15\linewidth}@{}p{.15\linewidth}p{.15\linewidth}p{.15\linewidth}p{.15\linewidth}p{.15\linewidth}}
    \# & \multicolumn{2}{c}{Player} && \multicolumn{2}{c}{Game Master} \\
    \hline

    \theutterance \stepcounter{utterance}  
    & & & \multicolumn{4}{p{0.6\linewidth}}{
        \cellcolor[rgb]{0.9,0.9,0.9}{
            \makecell[{{p{\linewidth}}}]{
                \texttt{\tiny{[P1$\langle$GM]}}
                \texttt{You are an expert in causal inference. The following question is not a typical commonsense query, but rather a meticulously designed question created by a professor specializing in causal inference, intended to assess the students' mastery of the course content.} \\
\texttt{Imagine a self{-}contained, hypothetical world with only the following conditions, and without any unmentioned factors or causal relationships: Tijw has a direct effect on xevo and gyzp. Xevo has a direct effect on tijv. Tijv has a direct effect on gyzp. Tijw is unobserved. We know that tijw causes xevo. xevo causes tijv. tijw or tijv causes gyzp. We observed an individual is not tijw. Would an individual is not gyzp if not xevo instead of xevo?} \\
\texttt{Start your answer with 'Yes' or 'No', followed by additional reasoning or evidence to support your explanation.} \\
            }
        }
    }
     \\ \\

    \theutterance \stepcounter{utterance}  
    & \multicolumn{4}{p{0.6\linewidth}}{
        \cellcolor[rgb]{0.561,0.737,0.561}{
            \makecell[{{p{\linewidth}}}]{
                \texttt{\tiny{[P1$\rangle$GM]}}
                \texttt{No.} \\
\\ 
\texttt{To address this question, we need to apply the principles of causal inference, particularly the concept of causal graphs and the backdoor criterion. } \\
\\ 
\texttt{Given the causal relationships:} \\
\\ 
\texttt{{-} Tijw causes xevo} \\
\texttt{{-} Xevo causes tijv} \\
\texttt{{-} Tijw or tijv causes gyzp} \\
\\ 
\texttt{We can represent these relationships as a causal graph:} \\
\\ 
\texttt{Tijw {-}> Xevo} \\
\texttt{Xevo {-}> Tijv} \\
\texttt{Tijw {-}> Gyzp} \\
\texttt{Tijv {-}> Gyzp} \\
\\ 
\texttt{Since Tijw is unobserved, we cannot directly observe its effect on Gyzp. However, we know that Tijw causes Xevo, and Xevo causes Tijv, which in turn causes Gyzp. This creates a path from Tijw to Gyzp through Xevo and Tijv.} \\
\\ 
\texttt{Now, let's consider the scenario where the individual is not Tijw. We want to know whether the individual is not Gyzp if they were not Xevo instead of Xevo.} \\
\\ 
\texttt{To answer this question, we need to consider the backdoor criterion, which states that a path from a cause to an effect is blocked if we can control for all intermediate variables that are affected by the cause and affect the effect.} \\
\\ 
\texttt{In this case, the path from Tijw to Gyzp through Xevo and Tijv is blocked because we know that the individual is not Tijw. However, the path from Tijw to Gyzp through Xevo is not blocked because we are considering an alternative scenario where the individual is not Xevo.} \\
\\ 
\texttt{To block this path, we need to control for Tijv, which is affected by Xevo. However, we are not given any information about the effect of not Xevo on Tijv. Therefore, we cannot conclude that the individual is not Gyzp if they were not Xevo instead of Xevo.} \\
\\ 
\texttt{In other words, the absence of Xevo does not necessarily imply the absence of Gyzp, because Tijv may still be caused by Tijw, which is unobserved. Therefore, the answer is No.} \\
            }
        }
    }
    & & \\ \\

    \theutterance \stepcounter{utterance}  
    & & & \multicolumn{2}{p{0.3\linewidth}}{
        \cellcolor[rgb]{0.9,0.9,0.9}{
            \makecell[{{p{\linewidth}}}]{
                \texttt{\tiny{[GM$|$GM]}}
                \texttt{No} \\
            }
        }
    }
    & & \\ \\

    \theutterance \stepcounter{utterance}  
    & & & \multicolumn{2}{p{0.3\linewidth}}{
        \cellcolor[rgb]{0.9,0.9,0.9}{
            \makecell[{{p{\linewidth}}}]{
                \texttt{\tiny{[GM$|$GM]}}
                \texttt{yes} \\
            }
        }
    }
    & & \\ \\

    \theutterance \stepcounter{utterance}  
    & & & \multicolumn{2}{p{0.3\linewidth}}{
        \cellcolor[rgb]{0.9,0.9,0.9}{
            \makecell[{{p{\linewidth}}}]{
                \texttt{\tiny{[GM$|$GM]}}
                \texttt{game\_result = LOSE} \\
            }
        }
    }
    & & \\ \\

\end{supertabular}
}

\end{document}
